\documentclass[11pt]{article}
\usepackage[margin=1in]{geometry}
\usepackage{amsmath,amssymb}
\usepackage[T1]{fontenc}
\usepackage{lmodern}
\usepackage{microtype}
\title{OBGEL Model Sheet: Cu--Ag -- Liquid}
\author{Open Binary Gibbs Energy Library}
\date{v0.2.0 candidate}
\begin{document}
\maketitle
\section*{Phase}
The Liquid phase is represented as a substitutional solution model.

\section*{Composition}
The independent composition variable is $x=x_{\mathrm{Cu}}$ with $x_{\mathrm{Ag}}=1-x$.

\section*{Molar Gibbs free energy}
\begin{align*}
G_m(x,T) &= xG^0_{\mathrm{Cu}}(T)+(1-x)G^0_{\mathrm{Ag}}(T)\\
&\quad+RT\left[x\ln x+(1-x)\ln(1-x)\right]\\
&\quad+x(1-x)\sum_k L_k(T)(2x-1)^k.
\end{align*}

\section*{Redlich--Kister parameters}
\begin{align*}
L_0(T)&=17384.4 - 4.46438T,\\
L_1(T)&=-1660.74 + 2.31516T.
\end{align*}

\section*{Standard-state Gibbs energies}
The pure-component functions are piecewise functions from the SGTE data for pure elements. The complete OBGEL model is limited by the common documented range of these functions.

\subsection*{Cu in Liquid}
\begin{align*}
G^0_{Cu}(T) &= 5194.277 + 120.97333T - 24.112392T\ln T - 0.00265684T^{2}\\[4pt]
&\quad+ 1.29223\times 10^{-7}T^{3} + 52478T^{-1} - 5.849\times 10^{-21}T^{7} && 298.15 \le T \le 1357.77\\[4pt]
&\quad - 46.545 + 173.88148T - 31.38T\ln T && 1357.77 \le T \le 3200
\end{align*}

\subsection*{Ag in Liquid}
\begin{align*}
G^0_{Ag}(T) &= 3815.564 + 109.31099T - 23.846331T\ln T - 0.001790585T^{2}\\[4pt]
&\quad - 3.98587\times 10^{-7}T^{3} - 12011T^{-1}\\[4pt]
&\quad - 1.034\times 10^{-20}T^{7} && 298.15 \le T \le 1234.93\\[4pt]
&\quad - 3587.111 + 180.96466T - 33.472T\ln T && 1234.93 \le T \le 3000
\end{align*}

\section*{Source and validation}
The binary interaction parameters are reproduced from Kusoffsky (2002), which attributes them to a 1998 private communication from H. L. Lukas. NIST identifies the Hayes--Lukas--Effenberg--Petzow assessment as the full Cu--Ag thermodynamic assessment. Independent numerical evaluation of the complete Cu--Ag model gives a eutectic temperature of approximately 1053 K (779.9 C), consistent with the NIST value of 779.1 C.

\end{document}
